\section*{Limitations}
Applications of today's LLM-based chatbots are constantly expanding. This work focuses only on knowledge-intensive dialogue. As such, other settings where chatbots are used to perform tasks (e.g. ``write an email for me'') or perform personalized chitchat (e.g. companion chatbots) might not benefit from \system's pipeline. More targeted study and evaluation of these settings is required to properly balance the generative power of LLMs and the factual knowledge injected from an external corpus for these applications. Settings where a chatbot needs to exhibit initiatives (e.g. try to be persuasive) are also outside the scope of this paper.
We also only consider one-hop information retrieval, because it is the most natural and prevalent form of users' information need. Multi-hop retrieval could improve information access for more complex queries. 

This work has only been tested on open-domain conversations. Generalizability of \system to specialized domains like medical or legal has not been evaluated. 

This work has only been evaluated on English conversations. Extending this work to more languages will be limited by the quality of available LLMs and retrieval systems in those languages. Incorporating recent work on multi-lingual models~\cite{scao2022bloom} and information retrieval systems~\cite{colbertx2022} is a promising future direction.